\section{Conclusion}\label{sec:conclusion}
The main conclusion is that exact sparse network--simplex formulations are
controlled by unobserved circulation freedom. Classical disaggregation
already gives the hull; the observation-sensitive construction retains only
one coordinate per independent cycle of each unobserved block--state
subgraph. It therefore provides a direct hull description when those
subgraphs are forests and a smaller exact extension when cycles remain.
The count describes this construction and the associated completion by
individual products, rather than unrestricted extension complexity.

Eliminating the remaining state coordinates reveals a different structural
question: which coefficients are needed in the original coordinates?
Bounded block cycle rank permits finite support and recovery libraries and
a uniform rank-dependent coefficient bound. Long flat chains admit a second
positive result: unit flow/product coefficients suffice for every observation
pattern with at most three observed labels, and four labels can force a
ratio of two. The rank-three $K_4$ example, sparse Fibonacci constructions
on simple series--parallel graphs of maximum degree three, and two-label
universality transfer show why neither small treewidth alone nor a small
number of labels alone implies a uniform coefficient bound over general
networks. These are obstructions to simple projected descriptions; exact
extended formulations and polynomial-time separation remain available.

For modeling, the results give a choice between retaining a small set of
state coordinates and generating complete structural cuts. For exact
verification, the recovery constructions turn membership into explicit
state flows and convex decompositions, including on degenerate boundaries.
The computational evidence supports substantial formulation-size savings,
particularly when labels occur in different blocks. It also shows that
global merging of unused labels and native flow bounds often make a
conventional LP the fastest method. A runtime benefit must therefore be
demonstrated for the intended workload rather than inferred from a smaller
formulation or an operation bound.

The scope is the equality-flow component hull with a shared simplex and
selected products. Additional side constraints may be imposed on this hull
to obtain a valid relaxation, but exactness for a larger coupled model
requires a separate argument. Likewise, the sharp observed-label threshold
is proved for flat chains, not for all nested series--parallel networks.
These boundaries separate the completed structural results from extensions
that the present arguments do not establish.
